% !TEX root = ../../cookbook.tex
\newpage
\recipe{Meringues\label{meringues}}{Pour environ 20 meringues.}{}{}

\subsection*{Ingrédients}
\begin{ingredients}
	\item 4 blancs d'\oe uf
	\item 300 g de sucre semoule
\end{ingredients}

\medskip

% --- Preparation ---
Préchauffer le four é 120°C, et préparer les \textbf{blancs d'\oe uf} à température ambiante. 
Commencer par battre les blancs à vitesse moyenne-douce. 
Quand ils commencent à mousser, s'il est disponible, ajouter quelques gouttes d'acide (jus de citron ou vinaigre blanc\footnote{L'acide aide à la rétention des bulles d'air, et facilite donc l'obtention d'un blanc bien ferme. Pour les géologues, l'HCl à 10\% n'est pas une alternative recommandée.}).

Dès que les fouets commencent à laisser des traces dans la mousse, augmenter la vitesse à moyenne-forte et ajouter \sfrac{1}{3} du \textbf{sucre} en pluie.

Attendre 1 ou 2 minutes que le sucre soit bien incorporé, puis ajouter le deuxième tiers.

Ajouter le dernier tiers lorsque la meringue devient bien blanche et commence à prendre du volume.

Serrer la méringue à vitesse maximale jusqu'à ce qu'elle forme un \textit{bec d'oiseau} bien ferme.

Dresser\footnote{Il n'est pas nécessaire de les pocher, il suffit de faire des petites formes avec une cuillère.} les meringues sur une plaque de four recouverte de papier sulfurisé, puis cuire pendant 20 minutes à 120°C, puis 1 h à 100°C. 


\vspace*{0.75 cm}

% --- Description ---
\begin{recipeblock}
	\noindent 
\end{recipeblock}